\documentclass[11pt]{article}
\usepackage[T1]{fontenc}
\usepackage{lmodern,microtype,amsmath,amssymb,amsthm,mathtools,booktabs,array,xcolor,enumitem}
\usepackage[a4paper,margin=26mm,top=24mm,bottom=25mm]{geometry}
\usepackage{fancyhdr}
\usepackage[colorlinks=true,linkcolor=blue!35!black,urlcolor=teal!55!black,citecolor=teal!65!black]{hyperref}
\definecolor{ink}{HTML}{143C48}
\setlist{itemsep=2pt,topsep=4pt}
\setlength{\headheight}{15pt}
\pagestyle{fancy}\fancyhf{}\fancyhead[L]{\small\color{ink}Sparse equations for singular moduli}\fancyhead[R]{\small\color{ink}Research note}\fancyfoot[C]{\small\thepage}
\newtheorem{theorem}{Theorem}[section]
\newtheorem{lemma}[theorem]{Lemma}
\newtheorem{corollary}[theorem]{Corollary}
\newtheorem{proposition}[theorem]{Proposition}
\newcommand{\Q}{\mathbb Q}\newcommand{\Z}{\mathbb Z}\newcommand{\F}{\mathbb F}
\newcommand{\vp}{v_p}
\newcommand{\source}[2]{\href{#1}{#2}}
\setlength{\emergencystretch}{2em}\raggedbottom\hypersetup{pdftitle={Five-term equations for singular moduli}}
\begin{document}
\thispagestyle{plain}
\noindent{\color{ink}\sffamily\small RESEARCH NOTE\hfill 25 NOVEMBER 2026}
\vspace{14pt}
\begin{center}{\LARGE\bfseries\color{ink} Five-term equations for singular moduli}\\[7pt]
{\large An exact classification and lower bounds in arbitrary degree}
\end{center}
\vspace{3pt}
\noindent\textbf{Abstract.}
A singular modulus satisfies a nonzero rational polynomial with at most five terms if and only if its degree is at most four. The exponents are unrestricted. The proof is independent of RH and GRH. It combines local root bounds, separated complex conjugates and finite-field tests which retain information when exponent residues coincide. Every quintic singular modulus requires exactly six terms. In arbitrary degree, the Dirichlet theorem in \emph{grh-lean} gives the uniform estimate $s(\Delta)\gg h(\Delta)/(\log|\Delta|)^2$ for the least support. We also construct infinitely many explicit orders with $s=h+1$ and prove that full support holds for more than $61.33\%$ of Gaussian prime conductors in lower relative density. A separate application of the repository's one-sided argument gives a sharp character-theoretic optimization and isolates an open estimate for the splitting field of $X^4-2$.

\section{What the results say}
Write $\Delta=f^2d_K<0$ for the discriminant of an imaginary quadratic order of conductor $f$ in $K$. Its singular moduli are the $h=h(\Delta)$ distinct conjugates of $j(\tau)$, one for each reduced primitive positive quadratic form $[a,b,c]$ of discriminant $\Delta$, where $\tau=(-b+i\sqrt{|\Delta|})/(2a)$. Write
\[
 s(\Delta)=\min_{0\ne P\in\Q[X],\ P(j(\tau))=0}\#\{i:[X^i]P\ne0\}.
\]
Equivalently, minimize the number of terms of any nonzero multiple of the monic ring class polynomial $H_\Delta\in\Z[X]$. Degrees and coefficient sizes of multiples are unrestricted. Its own degree $h$ gives $s(\Delta)\le h+1$.

\begin{theorem}[Five-term classification]\label{thm:four}
For every imaginary quadratic order discriminant $\Delta$,
\[
       s(\Delta)\le5\quad\Longleftrightarrow\quad h(\Delta)\le4.
\]
Furthermore, $s(\Delta)=5$ whenever $h=4$, and $s(\Delta)=6$ whenever $h=5$. All three assertions hold independently of RH and GRH.
\end{theorem}
The easy direction uses the class polynomial itself. The content of the reverse direction is that for every singular modulus of degree at least five, \emph{any five distinct nonnegative powers} are linearly independent over $\Q$. Sections~\ref{sec:finite} and \ref{sec:quintic-sieve} give the exact small-degree minima. Sections~\ref{sec:global-five} and \ref{sec:four-sieves} exclude the remaining degrees. Complex multiplication also gives $[K(j):K]=[\Q(j):\Q]=h$. The rational basis $1,j,\ldots,j^{h-1}$ remains a basis over $K$, so every fixed set of powers has the same rank over the two fields. The support results extend to coefficients in the underlying imaginary quadratic field.

\paragraph{The questions in the preceding reports.}
Bilu, Luca and Pizarro-Madariaga treated trinomials. The ProveIt CM report treated four terms in degrees at least six and left 79 quartic and quintic discriminants. A companion report, \emph{Sparse equations for singular moduli and the sign of weighted prime sums}, gave a one-prime proof of the trinomial conjecture, proved the four-term classification independently of GRH, and established the six-term minimum for 28 of the 29 quintic orders. At the last quintic discriminant, $-131$, it excluded five terms through polynomial degree 1000. It asked whether five terms can occur in degree five or higher and whether its lower bound of order $\sqrt h$ can be improved. A second parallel report, \emph{Sparse equations for singular moduli}, reduced the global five-term question to precisely $-87,-131,-348,-432,-675$ and eliminated all five through exponent width $65{,}536$. It also derived the near-degree growth bound of Theorem~\ref{thm:growth} and an infinite family with $s=h+1$. Theorem~\ref{thm:four} closes all five remaining cases for arbitrary exponent width. Section~\ref{sec:family} gives a different exact family and a positive lower density among prime conductors.\cite{riffaut,prior,companion,sparseparallel}

\section{What transfers from the Lean proof}
The pinned source is \texttt{lsernau/grh-lean}, commit
\href{https://github.com/lsernau/grh-lean/commit/72fead102f663d5c7a407964a13dc0cf9d977207}{\texttt{72fead102f66}} (full identifier in the source audit). In the readable edition, \texttt{grh/src/Statements.lean} imports Mathlib and defines the public statements. The nonvanishing statement concerns every complex Dirichlet character $\chi$ of positive modulus, and every $s$ with $\Re s>1/2$, with the principal-character pole at $s=1$ excepted. The public declarations in \texttt{Main.lean} also give the zeta statement and the critical-strip statements for general and primitive characters. The stated toolchain is Lean 4.35.0 with Mathlib at \texttt{56a00585}. Andrew Sutherland independently records kernel replays and checks of the public statements for both editions. We inspected the declarations and matched thirteen pinned source files. Our new arguments are paper proofs with computational certificates.\cite{repo,mitcheck}

Two different uses of the repository matter here. The equality $\zeta_K=\zeta L(s,\chi_{d_K})$ supplies precisely the quadratic-field zero theorem required for uniform counts of split or inert primes. This is the input in the earlier large-discriminant cutoff and our global lower bound. The new local argument studies the reductions of the algebraic numbers themselves. At a suitable inert prime, many different conjugates have the same nonzero residue. A polynomial with few terms can have only a limited number of roots in such a disk, even if its degree is enormous.

The second use concerns the proof itself. The source compares an upper bound for a covolume derived from short, spanning forms with a lower bound from a determinant of theta lifts. The comparison yields a one-sided bound for the centered weighted sum built from
\[
 \sum_{\substack{n\le x\\n\equiv 1\bmod q}}\Lambda(n)
 \bigl(1-(n/d)\bigr)\bigl(1/n-1/x\bigr).
\]
The construction uses $18480\mid q$ and a prime $d\equiv q/2-1\pmod q$. A quadratic twist is chosen to avoid zeros in a prescribed box, which controls the signs of real residues. A Landau argument then detects zeros from the one-sided bound. Section~\ref{sec:transfer} uses this sign principle to ask which splitting statistics can see all characters beyond a known quotient. The arithmetic covolume bound for the Dirichlet problem has no corresponding specialization in the repository for our nonabelian example. Source paths and declaration names appear in \texttt{checks/source\_audit.md}.

\section{Counting sparse roots in one multiplicative disk}\label{sec:local}
Normalize $v_p(p)=1$ on $\overline{\Q}_p$. In the following arguments a finite complete extension contains all points and all scaling constants used. Write $\mathcal O$ for its valuation ring. Define, for a rational $\lambda>1/(p-1)$,
\[
 w_{p,\lambda}(n)=n\lambda-v_p(n!),\qquad w_{p,\lambda}(0)=0.
\]
\begin{lemma}[Local bound]\label{lem:disk}
Suppose $P$ has $k$ nonzero terms, $u\ne0$, and $v_p(b)=\lambda$. The number of distinct roots of $P$ in $u(1+b\mathcal O)$ is at most the largest integer $n$ such that
\[
 w_{p,\lambda}(n)\le\max_{0\le i<k}w_{p,\lambda}(i).
\]
The same bound holds on the corresponding open disk.
\end{lemma}
\begin{proof}
Use the bijection $X=u\exp(bY)$, $Y\in\mathcal O$, whose inverse is $\log(X/u)/b$. For $P(X)=\sum_{j=1}^k c_jX^{e_j}$ with distinct nonnegative integer exponents put $A_m=\sum_j c_j u^{e_j}e_j^m$, using $e_j^0=1$. Then
\[
 P(u\exp(bY))=\sum_{m\ge0} A_m\frac{b^m}{m!}Y^m.
\]
The first $k$ moments are not all zero by the ordinary Vandermonde determinant in the $e_j$. The monic polynomial $\prod_j(T-e_j)$ gives the sequence an integral linear recurrence. The ultrametric inequality therefore gives $|A_m|_p\le M=\max_{i<k}|A_i|_p$ for every $m$. At an index $i<k$ realizing $M$, the absolute coefficient is $M p^{-w_{p,\lambda}(i)}$. Every index past the displayed bound has strictly smaller coefficient. The series converges on the closed disk because $w_{p,\lambda}(m)\to\infty$. Strassmann's theorem bounds the number of zeros by the last index attaining the maximum coefficient norm. This proves the assertion. For rational $\lambda$, a finite extension realizing the valuation of $b$ always exists.
\end{proof}

This argument is related to Lenstra's local fewnomial bounds and to Avenda\~no and Krick's Proposition 3.8. Their proposition supplies the bound $k-1$ in a principal-unit disk if $p>e+k-1$, where $e$ is the ramification index. Retaining the individual factorial valuations in Lemma~\ref{lem:disk} handles small residue characteristic and some large values of $k$.\cite{aven}

\begin{corollary}\label{cor:coarse}
In an unramified local field at an odd $p$, a disk specified by one nonzero residue contains at most $2(k-1)$ distinct roots of a $k$-term polynomial. For fixed $p$ it contains at most $k-1+O_p(\log(2k))$.
\end{corollary}
\begin{proof}
Lift the residue to a unit $u$ and take $b=p$. Write $t=k-1$. Since $v_p(n!)\le n/(p-1)$, we have $w_{p,1}(n)\ge n(p-2)/(p-1)$ and $\max_{i\le t}w_{p,1}(i)\le t$. This proves the first estimate. More precisely, writing $s_p(n)$ for the sum of the base-$p$ digits gives
\[
 w_{p,1}(n)=\frac{(p-2)n+s_p(n)}{p-1}.
\]
For $i\le t$ the digit sum is at most $(p-1)(1+\log_p(\max(1,t)))$. Substituting in the lemma gives the second estimate.
\end{proof}

\paragraph{Counting the conjugates in a certified disk.}
If $H$ is monic and separable, compute $B_i$ by expanding $H(p^a(u+Y))=\sum_i B_iY^i$ for a rational integer $u$ prime to $p$. The smallest index minimizing $v_p(B_i)+i\lambda$ equals the number of roots for which $v_p(p^{-a}X-u)>\lambda$. This is the Newton polygon root-counting rule. To see the endpoint convention, substitute $Y=bZ$, divide by a coefficient of minimum valuation, and reduce. A root strictly inside contributes a factor $Z$, a root on the boundary a factor with nonzero constant term, and an exterior root contributes a constant after normalization. Thus the order at $Z=0$ of the reduction counts precisely the interior roots. Division of the center by its unit part converts the disk to that in Lemma~\ref{lem:disk}. Rational noninteger radii require no computations in a ramified field in the verifier.

\section{Small-degree certificates and the previous four-term problem}\label{sec:finite}
The earlier CM report used separated complex conjugates and a prime-ideal estimate to reduce the four-term problem to 79 order discriminants. We reproduced its large and small range sieves, the 121 exceptional minors at $\Delta=-4075$ and the constants in the analytic cutoff. The 79 cases comprise 50 quartic and all 29 quintic orders. Their independent certification is also a direct way to establish the small-degree inputs required for the five-term theorem. The earlier four-term reduction is documented in Appendix~\ref{sec:earlier} and its checkers under \texttt{previous/}.\cite{prior}
\paragraph{The final 79.}
For each remaining discriminant we computed $H_\Delta$ by enumerating its primitive reduced forms and evaluating
\[
 j(\tau)=\frac{E_4(q)^3}{q\prod_{n\ge1}(1-q^n)^{24}},\qquad
 E_4(q)=1+240\sum_{n\ge1}\sigma_3(n)q^n,\quad q=e^{2\pi i\tau}.
\]
All evaluations use outward-rounded dyadic intervals. Reduction gives $|q|<1/200$. The 164-term truncation of $E_4$ has error at most $960\cdot165^3/200^{165}$, since $\sigma_3(n)<2n^3$. For the omitted factors of the product the sum of $|q|^n$ is at most $2/200^{165}$, which gives an absolute full-product error below $8/200^{165}$. Each coefficient enclosure contains exactly one integer, determining the monic class polynomial by CM integrality. The checker independently reconstructs all 79 polynomials.

For each one, a row of \texttt{cm/disk\_witnesses.csv} specifies $p,a,u,\lambda$ in the Newton polygon test of Section~\ref{sec:local}. It verifies that at least four conjugates lie in the disk and that $w_{p,\lambda}(n)>\max_{i<4} w_{p,\lambda}(i)$ for all $n\ge4$. To test this infinite condition it uses exact factorial valuations up to the point where $n\lambda-(n-1)/(p-1)$ exceeds the required maximum and increases thereafter. Thus each certificate excludes all rational four-term relations, including those of arbitrarily large degree. The following counts describe the 79 four-term exclusions and the quintic subset for which one disk also excludes five terms.
\begin{center}
\begin{tabular}{lrrrrrr}
\toprule Witness prime&2&3&7&13&17&Total\\\midrule
Four-term exclusions&1&48&23&5&2&79\\
Five-term exclusions by a single disk&0&20&5&3&0&28\\\bottomrule
\end{tabular}
\end{center}

\paragraph{One certificate by hand.}
At $\Delta=-68$ expand $2^{-24}H_{-68}(2^6(3+Y))$. From constant to leading term the valuations are $13,10,6,5,0$. With $\lambda=5/2$ the unique minimizing index of $v_2(B_i)+5i/2$ is 4. All four conjugates are inside the disk. For $i=0,1,2,3$, the maximum of $5i/2-v_2(i!)$ is $13/2$. At $n=4$ the weight is $7$, and for $n\ge5$ it is at least $3n/2+1$. Strassmann gives at most three sparse roots for any polynomial of support at most four. This proves $s(-68)=5$ with a small integer calculation.

\paragraph{An independent census in degrees four and five.}
These two degrees admit a proof independent of GRH. The maximal order class number $h_K$ divides the order class number. For $h=4$ the possible $h_K$ are 1, 2 and 4, and Watkins's Table 4 gives upper bounds $163,427,1555$, respectively, for the absolute field discriminant. For $h=5$ the options are 1 and 5 with bounds 163 and 2683. The conductor formula is
\[
 h=h_K f\prod_{p\mid f}(1-\chi_{d_K}(p)/p)/[\mathcal O_K^\times:\mathcal O_f^\times].
\]
The index is at most three, so $\varphi(f)\le 3h$. On prime powers the product formula gives $\varphi(f)^2\ge f/2$ (only $2^1$ can have $\varphi(p^a)^2<p^a$). Hence $f\le288$ or $450$, respectively. Enumerating the fundamental discriminants, reduced forms, and all these conductors gives 84 quartic and 29 quintic orders. The 34 quartic orders outside the 50 disk cases all pass the four-radius permutation bound with certified intervals. For 28 of the 29 quintic orders the disk check itself gives $s=6$. At $\Delta=-131$ the reduction modulo 17 is $X(X-8)^4$. The disk at 8 captures four of the five roots. The next section resolves the five-term question in this case. All checks in this paragraph are included under \texttt{cm/}.\cite{watkins}


\section{Five terms at discriminant \texorpdfstring{$-131$}{-131}}\label{sec:quintic-sieve}

The class polynomial for $\Delta=-131$ has degree five. We show that
\emph{any} five distinct nonnegative powers of a root are linearly
independent over $\Q$. The test uses six primes. Each prime restricts the
possible exponent residues, including the case when two exponents have the
same residue.

\begin{lemma}[A finite test for unbounded exponents]\label{lem:confluent}
Let $H\in\Z[X]$ be monic and irreducible, let $x$ be the class of $X$
in $A=\Z_{(p)}[X]/(H)$, and suppose $x^m=a+pD$ in $A$ for a
rational integer $a$ prime to $p$. Fix distinct nonnegative integers
$e_1,\ldots,e_k$ and let $n_r$ count those congruent to $r$ modulo $m$.
If $p>k+1$ and the following columns are linearly independent in
$A/pA$ over $\mathbb F_p$, then $x^{e_1},\ldots,x^{e_k}$ are linearly
independent over $\Q$:
\[
                x^r D^j,\qquad 0\le r<m,\quad 0\le j<n_r .
\]
\end{lemma}
\begin{proof}
Within each residue block write $e=r+mt$ and divide the column by $a^t$.
It becomes $x^r(1+pD/a)^t$. For the first $j+1$ distinct values of $t$
in this block, form the Newton divided difference of order $j$ and
multiply by $j!/p^j$. Doing this successively gives an invertible
triangular change of columns over $\Q$.

Expand by the binomial theorem up to the largest $t$ in the block.
A divided difference of order $j$ annihilates $\binom{T}{i}$ for $i<j$
and equals $1/j!$ for $i=j$. For $i>j$ its denominator divides $i!$:
the polynomial $i!\binom{T}{i}$ has integer coefficients, and divided
differences of integer polynomials at integer arguments are integers.
The latter identity follows on monomials, whose order-$j$ divided
differences are the complete homogeneous symmetric polynomials of
degree $i-j$ at the chosen arguments.
For $i>j$ we have $i-j>v_p(i!)$. This is immediate for $i<p$.
For $i\ge p$ use $j\le k-1\le p-3$ and
$v_p(i!)\le (i-1)/(p-1)<i-j$.
Every transformed column therefore belongs to $A$, with reduction
$x^r(D/a)^j$ modulo $p$. Its scalar $a^{-j}$ is nonzero. Independence
after reduction implies independence over $\Q$ since $A$ is free over
$\Z_{(p)}$ on $1,x,\ldots,x^{\deg H-1}$. The triangular column operations
then give the conclusion.
\end{proof}

\begin{theorem}\label{thm:131}
The least support for discriminant $-131$ is six.
\end{theorem}
\begin{proof}
Use the certified class polynomial $H=H_{-131}$ in the accompanying data.
For each pair $(p,m)$ below, the checker reduces $x^m$ modulo
$(H,p^2)$. The constant coefficient is a $p$-adic unit, and every
other coefficient is divisible by $p$. Taking $a$ equal to the constant
coefficient determines $D\bmod p$.

Suppose five powers are dependent. Since $x\ne0$, divide by the lowest
power to obtain the exponent $0$ and four other exponents. At any
modulus $N$ we retain the sorted quadruple of the four residues, allowing
repetitions and zeros. On lifting $n\mid N$, add arbitrary multiples of
$n$ to the four old residues to obtain all representatives below $N$.
For equal old residues it suffices to take nondecreasing lifts among
those positions. Sorting the resulting quadruple gives every possible
lift exactly once. At each indicated prime, keep a quadruple only if the
five columns of Lemma~\ref{lem:confluent} are dependent. The distinguished
exponent $0$ supplies the first column $1$, so this test is a
four-by-four determinant after deleting its pivot row.
\begin{center}\small
\begin{tabular}{rrrr}
\toprule
$N$ & $(p,m)$ tested at this stage & Quadruples tested & Remaining\\
\midrule
200    & $(601,200)$                    &68,685,050 &119,264\\
1,400  & $(1051,350),\ (4201,1400)$      &279,218,674&196\\
4,200  & $(6301,1050)$                  &15,876     &5,876\\
12,600 & $(757,252)$                    &475,956    &214,120\\
25,200 & $(337,112)$                    &3,425,920  &0\\
\bottomrule
\end{tabular}
\end{center}
The modulus at every stage is divisible by each period used there.
Thus a genuine exponent set would give a surviving quadruple at every
stage. The last count proves independence for all five-element sets.
A relation on fewer powers would give a dependent five-element set by
adjoining distinct powers. Finally $H$ itself has at most six terms.
The polynomial reduction and determinant checks are reproduced by
\texttt{python3 checks/verify\_last\_quintic.py} from the archive root.
A separately written basis recurrence, cofactor computation and lift enumerator reproduce all five rows in \texttt{checks/independent\_quintic.py}.
\end{proof}

\paragraph{Exact minima through degree five.}
Combining the two finite censuses proves $s=h+1$ for every order of class number four or five. The same equality holds for $1\le h\le3$ except at $j=0$, where $s=1$. A nonzero rational $j$ has a binomial linear polynomial. For $h=2$ one has $D\ge15$, and the principal conjugate has larger absolute value than the other, since $e^{\pi\sqrt D}-2079>e^{\pi\sqrt D/2}+2079$. A binomial has all its nonzero complex roots on one circle. For $h=3$ the result is Theorem 6.6 of \cite{riffaut}, independent of GRH. Our checker \texttt{cm/verify\_all\_cubics.py} also certifies all 25 cubic orders using one disk per order, in characteristics $2,3,7,13,17$. It uses the degree-three instance of the conductor census above with $|d_K|\le907$, $f\le162$. Thus the complete table through $h=5$ is independent of hypotheses on zeros.

\section{Reducing the global five-term question to 71 discriminants}\label{sec:global-five}
Here $\Delta=-D=f^2d_K$ and $h=h(\Delta)$. We exclude any rational
annihilator with at most five nonzero terms if $h\ge5$, apart from the
explicit residual list treated separately. All reductions in this
section are unconditional. They use the local root bound in Section~\ref{sec:local} and determinants built from five conjugates.

\paragraph{The local choices.}
At $p=7,13$ the unique supersingular $j$ is nonzero. If $p\nmid f$
and $p$ does not split in $K$, the ring class field has absolute
ramification index at most two above $p$. Deuring reduction therefore
puts all $h$ roots into one nonzero multiplicative disk of depth
at least $1/2$. Five distinct roots in this disk are impossible for a
five-term polynomial. For completeness, in the local bound take
$\lambda=1/2$ and $k=5$. For $n=5,6$ one has
$n/2-v_p(n!)>2$, and for $n\ge7$ one has
$n/2-v_p(n!)\ge n/2-(n-1)/6>2$.
The first five weights are at most $2$.

If $7\cdot13\mid f$ we may also use $p=19$ when $19\nmid f$ and $19$
is nonsplit. The order class number formula gives
$h\ge\varphi(f)/3\ge\varphi(91)/3=24$.
There are only two supersingular $j$-values in characteristic 19, both
nonzero. For instance, the coefficient of $x^{18}$ in
$(x^3+Ax+B)^9$ is $B(10A^3+8B^2)$ modulo 19. Vanishing gives the
values $18$ and $7$ for nonsingular curves. One disk therefore holds
at least twelve roots. Absolute ramification index at most two again
suffices.

\paragraph{Archimedean certificates from split and conductor primes.}
A split odd prime $p\nmid f$ permits factors $p^j$ in the leading
coefficient of a primitive reduced form. One can combine these
factors with the square of a squarefree divisor of the conductor.
Specifically take $a=t^2s$, where $t\mid f$ is odd and squarefree,
$s$ is odd and coprime to $t$, and each prime dividing $s$ splits and
does not divide $D$. Assume $D\ge4a^2$. Put $D'=D/t^2$ and choose
$v$ with its parity prescribed by $D'$, so $4\mid v^2+D'$.
Modulo each prime power dividing $s$, lift a unit square root of $-D'$
by Hensel's lemma. For each $q\mid t$ avoid the at most two roots
of $v^2+D'\equiv0\pmod q$. The Chinese remainder theorem gives a
representative with $|v|\le ts$. Taking $b=tv$ and
$c=(v^2+D')/(4s)$ gives $|b|\le a\le c$. The choices at primes dividing
$t$ ensure $q\nmid c$, and those at primes dividing $s$ ensure
$q\nmid b$. Thus the form $[a,b,c]$ is primitive and reduced,
after normalizing boundary signs if necessary.

For two available coefficients $a<b$ the sufficient integral checks
are $D\ge16a^2$ and
$9D(b-a)^2\ge16a^2b^2$. They give a factor exceeding 16 between
the corresponding moduli: the cusp estimate
$||j_a|-\exp(\pi\sqrt D/a)|\le2079$ gives
$U\ge e^{4\pi}>34\cdot2079$ and $U/V\ge e^{4\pi/3}>32$ for the two
exponentials. Four successive gaps give five rows. To see that these rows work for all
increasing integer exponents, order the radii increasingly and scale
successive columns by factors in ratios
$(r_jr_{j+1})^{-(e_{j+1}-e_j)/2}$. In any row, moving by one entry
away from the diagonal loses a factor at least four. The off-diagonal
sum is at most $2\sum_{m\ge1}4^{-m}=2/3$ of the diagonal magnitude.
The generalized Vandermonde determinant is nonzero.

The following cases are disjoint. In the last two rows both 7 and 13
divide $f$. If a local choice above applies, that settles the case.
Every available local prime in the row otherwise splits, and the
displayed leading coefficients then exist and pass the separation
checks at and above the indicated bound. The conductor-only last row
needs no splitting assumption.
\begin{center}\small
\begin{tabular}{lrr}
\toprule
Conductor condition & Five leading coefficients & $D$ sufficient\\
\midrule
$7,13\nmid f$ (both absent) & $1,7,13,49,91$ & 38,416\\
$7\mid f$, $13\nmid f$ & $1,13,49,169,637$ & 1,623,076\\
$13\mid f$, $7\nmid f$ & $1,7,49,169,343$ & 470,596\\
$7\cdot13\mid f$, $19\nmid f$ & $1,19,49,169,361$ & 521,284\\
$7\cdot13\cdot19\mid f$ & $1,49,169,361,8281$ & 274,299,844\\
\bottomrule
\end{tabular}
\end{center}

\paragraph{Checking the finite ranges.}
For each odd $p$, one has $p\mid f$ exactly when $p^2\mid D$.
Enumerating the values $D\equiv0,3\pmod4$ in the five ranges, with
these divisibility restrictions, gives the following counts.
The leading-coefficient test searches integers $a\le251$ using only
the exact inequalities above and a directly verified primitive
reduced form at each $a$.
\begin{center}\small
\begin{tabular}{rrrrr}
\toprule
Row & Candidates & Local & Five leading coefficients & Remaining\\
\midrule
1 &18,705&14,908&3,667&130\\
2 &16,464& 8,820&7,634&10\\
3 & 1,364&   768&  594&2\\
4 &    30&    16&   14&0\\
5 &    45&     0&   45&0\\
\midrule
Total &36,608&24,512&11,954&142\\
\bottomrule
\end{tabular}
\end{center}
Counting all primitive reduced forms for the 142 remaining values
gives 21 with $h<5$ and four with $h=5$ already covered by the
degree-five theorem. For 46 more, a prime
$p\in\{7,13,19,31,37,43\}$ satisfies
$(\Delta/p)=-1$ and $h>4(\lfloor p/12\rfloor+
\mathbf1_{p\equiv7\bmod12})$. All these primes are $1\pmod3$,
so $j=0$ is ordinary. The classical count of supersingular $j$-invariants
in these characteristics is the parenthesized number. One nonzero
residue disk must therefore hold at least five conjugates. The
remaining 71 discriminants have $87\le D\le4563$.
The complete integer sieve and its final list are reproduced by
\texttt{global\_five.py}.

\paragraph{Certifying 67 of the 71 cases.}
The five-conjugate complex determinant test works for 49 discriminants.
Here the successive radii can be closer than the factor 16 used in the global cutoff.
Let $0<r_1<\cdots<r_5$ be the radii of the selected conjugates.
For an off-diagonal permutation $\sigma$, compare its term to the diagonal term
in the determinant with exponents $e_1<\cdots<e_5$.
Writing $g_j=e_{j+1}-e_j$, this ratio equals
\[
 \prod_{j=1}^4
 \left(\frac{\prod_{\sigma(i)>j}r_i}{\prod_{i>j}r_i}\right)^{g_j}.
\]
Each numerator has the same number of factors as its denominator, which is the product of the largest radii of that cardinality. The ratio is largest when all gaps are one. The checker encloses the five radii using the dyadic evaluation of $j$ in Section~\ref{sec:finite}. With exponents $0,1,2,3,4$, the lower bound for the diagonal magnitude exceeds the sum of upper bounds for the 119 other terms. This proves independence for arbitrary exponents.

For 18 further discriminants, an exact Newton polygon count and Lemma~\ref{lem:disk} exclude five terms in a single disk. The remaining four have $D=87,348,432,675$, all of class number six. For the 18 disk tests and these four remaining cases, the checker derives the class polynomial from interval enclosures and CM integrality. The file \texttt{checks/five\_small\_witnesses.json} records the forms, coefficients and local parameters used, and \texttt{checks/verify\_five\_small.py} verifies that its inputs match the 71 discriminants output by the global reduction.

\section{The four residual sextic discriminants}\label{sec:four-sieves}

It remains to treat $D=87,348,432,675$, where $D=|\Delta|$.
Their degree-six class polynomials were computed by enclosing the product
over all six reduced forms. Apply Lemma~\ref{lem:confluent} to five powers
of a root $x$. After dividing by the lowest power, one exponent is zero.
In the basis $1,x,\ldots,x^5$ its column is $1$. Omitting the constant
coordinate therefore leaves a $5\times4$ matrix. At the prime $p$, a
residue with multiplicity $l$ among \emph{all five} exponents contributes
the confluent columns of orders $0$ through $l-1$, with the distinguished
column removed afterwards. A candidate survives exactly when every
four-by-four minor vanishes at every prime tested at the current stage.
This gives finite tests even when the integer exponents grow without bound
or two of them coincide modulo arbitrarily high prime powers.

The table records the complete enumeration. The period $m$ in each
pair is certified by reduction modulo $(H_{-D},p^2)$ as in
Theorem~\ref{thm:131}. Each step from $n$ to a multiple $N$ enumerates
all lifts of the preceding sorted quadruples. All four nonzero-position
residues may themselves be zero.
\begin{center}\footnotesize
\begin{tabular}{rrrrrr}
\toprule
$D$ & $n$ & $N$ & $(p,m)$ & Tested & Left\\
\midrule
87 & 1 & 102 & $(103,102)$ &4,780,230&13,125\\
 &102 &1,530 & $(1531,1530)$ &653,917,425&625\\
 &1,530&7,650& $(151,150)$ &390,625&915\\
 &18&486& $(487,486),(1459,486)$ &7,971,615&0\\
348&1&486& $(487,486),(1459,486)$ &2,353,329,126&0\\
432&1&36& $(109,36)$ &82,251&106\\
 &36&144& $(433,144)$ &26,212&0\\
675&1&40& $(241,40)$ &123,410&1\\
 &40&280& $(13,56)$ &2,401&0\\
\bottomrule
\end{tabular}
\end{center}
For $D=87$ the 915 survivors modulo $7650$ reduce to precisely 15 sorted
quadruples modulo $18=\gcd(7650,486)$. The last step for $87$ tests every
lift of those 15 possibilities to $486$. Any set of integer exponents would
give compatible residues on both sides, so the zero final count excludes
every such set. For the other discriminants the successive moduli already
divide one another. In all four cases any five powers are independent over
$\Q$.

Run \texttt{python3 checks/verify\_five\_sieves.py} from the archive root.
The program reads the certified integer class polynomial coefficients,
constructs every finite-field column, compiles the C++17 enumerator, and
checks each count. The enumerator streams the candidates and stores only
the surviving quadruples. Thus even the stage testing more than two
billion candidates needs room for only its residue tables and survivors.
The full direct replay took 33 seconds including compilation on the
verification machine. Its largest list had 13,125 quadruples, and
the largest coefficient table had $1530\cdot5\cdot5$ field entries.

A second implementation, \texttt{checks/fast\_sieve/verify.py}, checks $D=87,348$ by eliminating two prefix columns and grouping suffix columns by their projective classes in the quotient. For $D=87$, the $196{,}587$ simultaneous failures modulo 486, the $13{,}125$ failures modulo 102 and the $19{,}669$ failures modulo 150 restrict the lifts to 979,524 candidates at modulus 1530. The rank test at $1531$ eliminates all of them. For $D=348$ the independent calculation already finds zero modulo 486.

Combining these four exclusions with the preceding 67, the unbounded-exponent quintic test and the global cutoff proves the reverse implication in Theorem~\ref{thm:four}. For $h\le4$, the degree of $H_\Delta$ proves the forward implication.

\section{An infinite family attaining the exact maximum}\label{sec:family}
\begin{theorem}\label{thm:family}
For every integer $m\ge1$, take
\[
 \Delta_m=-2^{\,6\cdot7^{m-1}},\qquad h_m=2^{\,3\cdot7^{m-1}-2}.
\]
Then $h(\Delta_m)=h_m$ and $s(\Delta_m)=h_m+1$. This statement is independent of GRH.
\end{theorem}
The family in \cite{sparseparallel} has discriminants $-2^{2+6\cdot7^{m-1}}$ for $m\ge2$. The discriminants here differ by a factor of four. In particular, every choice of $h_m$ distinct nonnegative powers of such a $j$ is a basis for $\Q(j)$ over $\Q$. Dependence would give a polynomial with at most $h_m$ terms. Taking the exponents $0,1,\ldots,h_m$ shows that the minimum is attained. At $m=2$ it is $524{,}289$. At $m=3$ it is $2^{145}+1$.

\begin{proof}
Write $M=7^m$, $h=2^{3\cdot7^{m-1}-2}$, $f=2h$. The order of conductor $f$ in $\Q(i)$ has discriminant $-4f^2=\Delta_m$. Its class number is $f/2=h$, by the order class number formula: the Gaussian field has class number one, the unit index is two, and 2 ramifies.

The prime 7 is inert in $\Q(i)$ and coprime to $f$. The ring class field is unramified above 7. By Deuring reduction every conjugate of $j$ reduces to a supersingular invariant. In characteristic 7 the only such invariant is $1728\equiv6$. To check uniqueness directly, the Hasse invariant of $y^2=x^3+Ax+B$ in characteristic 7 is the coefficient of $x^6$ in $(x^3+Ax+B)^3$, namely $3B$. Supersingularity forces $B=0$. Thus all $h$ distinct conjugates lie in the same disk $6(1+7\mathcal O)$ in a common unramified completion.\cite{cmred}

Put $w(n)=n-v_7(n!)$. We show
\[
  w(i)<w(h)\ (0\le i<h),\qquad w(h)<w(i)\ (i>h).
\]
Because $2^3=1+7$, raising this equality to the power $7^{m-1}$ gives $4h\equiv1\pmod M$. For an interval of $r$ consecutive positive integers adjoining the cut, either $[h-r+1,h]$ (when $r\le h$) or $[h+1,h+r]$, it suffices to prove that the sum of the 7-adic valuations in the interval is strictly less than $r$.

First take $1\le r<M/5$. Its elements have the form $h+j$ with $|j|\le r$. The nonzero integer $1+4j$ has $|1+4j|\le4r+1<M$, so $v_7(h+j)=v_7(1+4j)$. In either arithmetic progression of $r$ consecutive choices of $j$, there are at most $r/7^a+1$ multiples of $7^a$ among the values $1+4j$. Summing over $a\le\log_7(4r+1)$ bounds the valuation sum by
\[
    r/6+\log_7(4r+1)<r .
\]
Indeed $4r+1\le5^r$ and $5^6<7^5$.

If $r\ge M/5$, direct counting of multiples in the original interval instead gives $r/6+\log_7(h+r)$. When $r\le h$, we have
\[
 \log_7(h+r)\le\log_7(2h)<M/6\le5r/6 ,
\]
where $2^{18}<7^7$ implies $(3/7)\log_7 2<1/6$. When $r>h$, use $h+r<2r\le 5^r$ and $\log_7 5<5/6$. These exhaust the two intervals and prove the strict inequalities for $w$.

It follows that $\min_{n\ge h} w(n)>\max_{0\le i<h}w(i)$. Lemma~\ref{lem:disk} prevents a polynomial with at most $h$ terms from vanishing at $h$ points of the disk. The degree-$h$ class polynomial supplies the opposite inequality.
\end{proof}

\paragraph{How frequently a single residue helps.}
For any order with $(\Delta/7)=-1$ or $(\Delta/13)=-1$, the same argument through Corollary~\ref{cor:coarse} yields
$s(\Delta)\ge h-O(\log(2h))$, independently of GRH. In characteristic 13 there is also exactly one supersingular invariant, and it is nonzero (the standard count is $\lfloor p/12\rfloor$ plus the corrections for 0 and 1728). The conductor is prime to an inert prime detected by the symbol. Among order discriminants counted once by absolute value, the stated union has natural density $1-(1-3/7)(1-6/13)=9/13$ by the Chinese remainder theorem. This is the density of the sufficient congruence conditions.
\subsection*{Full support for a positive proportion of prime conductors}
Let $\ell\ne2,7$ be prime, $\Delta=-4\ell^2$ and
$h=(\ell-\chi_{-4}(\ell))/2$. All $h$ conjugates reduce to the same
nonzero supersingular $j$ in characteristic $7$. Write
$w(n)=n-v_7(n!)$. By the local argument, $s(\Delta)=h+1$ whenever
\[
 w(u)<w(h)<w(v)\qquad (0\le u<h<v).
\]
Among primes $\ell$, this sufficient condition has a relative natural
density $d>0.6133$. Consequently the lower relative density of primes
giving $s(-4\ell^2)=h+1$ exceeds $61.33\%$. This result is independent
of Dirichlet GRH.

\smallskip\noindent\emph{Proof.}
An obstruction at distance $r$ requires
$\sum_{t\in I}v_7(t)\ge r$ on at least one of
$I=[h-r+1,h]$ (when $r\le h$) and $I=[h+1,h+r]$.
Let $b=\max_{t\in I}v_7(t)$. Counting multiples of $7^a$ in the
interval gives
\[
 \sum_{t\in I}v_7(t)\le\sum_{a=1}^b(r/7^a+1)\le r/6+b.
\]
Hence $q_r=7^{\lceil5r/6\rceil}$ divides $h+j$ for an offset
$1-r\le j\le r$. For fixed $r$, the prime number theorem in
progressions, applied separately to each sign $\sigma=\chi_{-4}(\ell)$
with $\ell=2h+\sigma$, bounds the relative density of the obstruction by
$u_r=2r/\varphi(q_r)$. Indeed there are at most $2r$ admissible classes
of $\ell$ per sign modulo $4q_r$, each of density $1/(2\varphi(q_r))$.

For completeness, these fixed-modulus bounds can be summed to control
arbitrarily distant obstructions. For primes $\ell\le X$ and
$q_r\le\sqrt X$, Brun--Titchmarsh bounds the count for each coprime
residue modulo $q_r$ by $2X/(\varphi(q_r)\log(X/q_r))$.
Ignoring the congruence modulo $4$, both signs and all offsets cost at
most $4r$ such classes. Thus for $r>R$ in this range the total is at
most $(X/\log X)\sum_{r>R}16r/\varphi(q_r)$. If $q_r>\sqrt X$,
a positive integer $h+j$ in the interval is divisible by
$7^{\lfloor\log_7\sqrt X\rfloor}$. Necessarily
$5r/6\le\log_7(h+r)$. For $r>h$ this fails, since $2r\le2^r$ and
$\log_7 2<5/6$. For $r\le h$ it implies $r=O(\log X)$.
There are $O(\sqrt X)$ possible multiples and $O(\log X)$ offsets,
hence $O(\sqrt X\log X)$ possible primes. This is $o(\pi(X))$.
The sieve inequality used here is due to Montgomery and Vaughan,
\emph{The large sieve}, Mathematika 20 (1973), 119--134
(\url{https://doi.org/10.1112/S0025579300004708}).

For fixed $R$, the event of no obstruction at distances $r\le R$ has
a density $d_R$, computable using complete residue classes modulo
$4\cdot7^R$: every valuation may be truncated at $R$ without affecting
a comparison with $r\le R$. The finitely many $h<R$ are immaterial.
The tail estimate just proved and convergence of $\sum u_r$ imply that
the density of no obstruction at any distance exists and equals
$d=\lim_Rd_R$. Applying the fixed-$r$ union bound first on
$R<r\le T$ and then letting $T\to\infty$ improves the error to
\[
                  0\le d_R-d\le\sum_{r>R}u_r.
\]
For instance, if $a=R+1$ and $b=(1+6/a)7^{-5}$, then
$u_{r+6}/u_r=(1+6/r)7^{-5}$ gives the explicit bound
$\sum_{r>R}u_r\le (u_a+\cdots+u_{a+5})/(1-b)$.

The supplied exact calculation gives
\[
 d_{10}=\frac{14143023}{23059204},\qquad
 d_{40}=\frac{33191837301111812268560181397}
              {54116956037952111668959660849},\qquad
 0\le d_{40}-d<5.6\cdot10^{-28}.
\]
To reproduce it, start with the twelve pairs consisting of a sign
$\sigma$ and a residue $h$ modulo $7$ with $7\nmid 2h+\sigma$.
Each pair has mass $1/12$. At depth $v$, compute each interval sum
with valuations capped at $v$. Discard the class if a sum already
reaches its interval length. Keep the class if every sum is smaller
and every valuation below $v$. Otherwise split into its seven lifts,
each with one seventh of the mass. At depth $R$, every ambiguous
interval already has truncated sum at least $R$ and is discarded.
Summing the kept masses yields $d_R$.
Thus $d$ rounds to $0.6133352599845868$ at sixteen places after the
decimal point. The exact-support set could contain further primes
for which this sufficient condition fails.\hfill$\square$


\section{A uniform lower bound for arbitrary orders}\label{sec:uniform}
\begin{theorem}\label{thm:growth}
With Dirichlet GRH as supplied by the repository, there is an absolute constant $c>0$ such that for all sufficiently large $D=|\Delta|$,
\[
 s(\Delta)\ \ge\ c\,\frac{h(\Delta)}{(\log D)^2}
 \ \gg\ \frac{\sqrt D}{(\log D)^2\log\log D}.
\]
Consequently $\log s(\Delta)/\log D\to1/2$ through all imaginary quadratic order discriminants.
\end{theorem}
The earlier ProveIt bound used many conjugates separated in ordinary absolute value and gave $D^{1/4}/\sqrt{\log D}$. The companion report subsequently gave $s\gg\sqrt h$ with Dirichlet GRH and $s\gg_\epsilon D^{1/4-\epsilon}$ independently.\cite{companion} The later parallel report independently gives both exponents stated in this section.\cite{sparseparallel} The exponent comes from counting conjugates with the same reduction at a small inert prime. The distribution of primes gives a supply of possible characteristics. A global norm argument then finds one in which most of the conjugates are units.

\begin{proof}
Write $h=h(\Delta)$ and $D=f^2|d_K|$. We record the uniform inputs explicitly.

\emph{Class number and norm.}
For the possibly imprimitive Kronecker character $\chi_\Delta$ of modulus $D$, the analytic class number formula for orders gives
$h=w_\Delta\sqrt D L(1,\chi_\Delta)/(2\pi)$, with $w_\Delta\in\{2,4,6\}$ the number of units. Under Dirichlet GRH one has $L(1,\chi_\Delta)\gg1/\log\log D$ uniformly in $f$. For detail, the standard square-root estimate for the primitive character follows by subtracting the prime ideal sums of $\Q$ from those of $K$. The primes removed on induction to modulus $D$ cost $O(\log(2D)\log(2t))$ in $\sum_{n\le t}\Lambda(n)\chi(n)$. Thus the total character sum with von Mangoldt weight is $O(\sqrt t\log^2(2Dt))$. By partial summation the tail of $\sum_{n}\Lambda(n)\chi_\Delta(n)/(n\log n)$ beyond $Y=(\log D)^8$ is $o(1)$ uniformly. This sum equals $\log L(1,\chi_\Delta)$ by passage from $s>1$. The partial sum is at least $-\sum_{p\le Y}1/p-O(1)$ since all higher powers contribute absolutely $O(1)$. Mertens' estimate gives the claimed bound after exponentiation. The same proof with the signs reversed gives $h\ll\sqrt D\log\log D$.

Put $N=|H_\Delta(0)|$. It is nonzero away from $\Delta=-3$. Sutherland's Lemma 8, including its proof for quadratic orders, gives
\[
 \log N\ll\sqrt D\log D\log\log D .
\]
The GRH ingredient of this height estimate is the quadratic-character bound of Littlewood, as identified in the proof of that lemma. Hence $\log N/h\ll\log D(\log\log D)^2$.\cite{suth,gm}

\emph{A supply of inert primes.}
Let $X=C(\log D)^2$ for a sufficiently large absolute $C$. For $M=\Q,K$ put
\[
 W_M(X)=\sum_{N\mathfrak a\le X}\Lambda_M(\mathfrak a)(1-N\mathfrak a/X).
\]
The smoothed prime ideal theorem, divided by $X$, gives
$W_M(X)=X/2+O(\sqrt X\log(2D)+\log X+\log(2D))$
in these degrees, with an absolute constant. Both zeta hypotheses follow from Dirichlet RH and GRH. In $2W_{\Q}-W_K$, first powers of split primes cancel, inert primes have coefficient two, and ramified primes coefficient one. Prime powers of rational base $p$ with exponent at least two cost $O(\sqrt X\log^2(2X))$ by elementary counting, and the ramified first powers and the primes dividing $f$ cost $O(\log D)$. Dropping primes below 5 is harmless. Comparing with the main term $X/2$, and using $0\le1-p/X\le1$, gives
\[
 \sum_{p\in\mathcal P}\log p\gg X,\qquad
 \mathcal P=\{5\le p\le X:\ p\text{ inert in }K,\ p\nmid f\}.
\]
The error $O(\sqrt X\log(2D))$ is a sufficiently small fraction of $X$ once $C$ is large.\cite{gms}

\emph{Selecting the disk.}
Choose a place of the ring class field over each $p\in\mathcal P$. This field is unramified locally, and all roots of $H_\Delta$ are integral. Each of the $z_p$ roots reducing to zero has normalized valuation at least one, so $z_p\le v_p(N)$. Consequently
\[
 \sum_{p\in\mathcal P} z_p\log p\le\log N
       =o(hX).
\]
For some $p\in\mathcal P$ at least $h/2$ of the roots are units. Every reduction is supersingular by Deuring's theorem. Over the algebraic closure of $\F_p$ there are $p/12+O(1)$ such invariants. At least $\gg h/p$ distinct roots therefore lie in one disk above a nonzero residue. Corollary~\ref{cor:coarse} gives $2(s(\Delta)-1)\gg h/p$, proving the lower bound. Combining it with $s(\Delta)\le h+1$ and the upper class number estimate proves the limit.\cite{cmred,supersing}
\end{proof}

\begin{corollary}[A comparison independent of GRH]\label{cor:pollack}
For every $\epsilon>0$,
\[
 s(\Delta)\gg_\epsilon D^{1/2-1/(4\sqrt{\mathrm e})-\epsilon}
\]
for all sufficiently large $D$. The constant obtained by this argument is ineffective.
\end{corollary}
\begin{proof}
Apply Pollack's Theorem 1.1 to the nontrivial quadratic character $\chi_\Delta$ modulo $D$. For any $\delta>0$ it yields $>D^\kappa$ prime nonresidues below $D^{1/(4\sqrt{\mathrm e})+\delta}$, with $\kappa=\kappa(\delta)>0$. They are inert and coprime to $D$. The unconditional half of Sutherland's Lemma 8 gives $\log N\ll\sqrt D\log^2D$, and Siegel's class number bound and the order formula give $h\gg_\eta D^{1/2-\eta}$ for every $\eta>0$. For the latter, use the maximal-order bound and $\prod_{p\mid f}(1-\chi_{d_K}(p)/p)\ge\varphi(f)/f\gg1/\log\log(3f)$, with unit index at most three. Choose $\delta<\epsilon/2$ and $\eta<\min(\epsilon/2,\kappa/2)$. The same norm average, whose denominator now exceeds a constant multiple of $D^\kappa$, gives a prime with at least $h/2$ nonzero reductions. The disk argument gives $s(\Delta)\gg h/D^{1/(4\sqrt{\mathrm e})+\delta}$. The use of Siegel's bound accounts for ineffectivity.\cite{pollack,suth}
\end{proof}

\section{How sparse can a one-sided Frobenius test be?}
\label{sec:transfer}
The residue argument in
\texttt{New.LFunction\_ne\_zero\_of\_primeRemainder\_le} suggests a
finite extremal question. Once a quotient of a Galois group has been
controlled, can a bound for a single nonidentity splitting type see
every remaining irreducible character with the required sign? We give
a sharp answer for tests designed by the signs of individual character
coefficients. This qualification matters because additional arithmetic
relations among zero multiplicities can allow other tests.

\subsection*{An exact support and minimax theorem}
Let $G$ be a finite group of order $g$, and let $N\ne1$ be a normal
subgroup of order $n$. Call a complex irreducible character
$\chi$ \emph{new} in this section if its representation is nontrivial
on $N$. All other irreducibles come from $G/N$. Let $f$ be a real
class function on $G$ with $f(u^{-1})=f(u)$, $0\le f\le 1$, and $f(1)=0$.
Write $f=\sum_\chi a_\chi\chi$, with the usual normalized character
inner product as coefficient. The $a_\chi$ are real by inversion
symmetry. Suppose $a_\chi\le0$ for every new $\chi$.
Define the worst response to one unit of regular-character multiplicity by
\[
  t(f)=\min_{\chi\ \mathrm{new}}\frac{-a_\chi}{\chi(1)}.
\]

\noindent\textbf{Theorem.}
The zero set of $f$ on $N$ is exactly
$N\cap\bigcap_{\chi\ \mathrm{new},\ a_\chi<0}\ker\chi$.
If every new coefficient is strictly negative, the support contains
$N\setminus\{1\}$. The indicator
$f_*=\mathbf 1_N-\mathbf 1_{\{1\}}$ attains this minimal support.
Moreover,
\[
                      \max_f t(f)=1/g.
\]
The maximizers are precisely $f_*+h$ for real inversion-symmetric
class functions $h$ inflated from $G/N$, with $h(N)=0$ on the identity
coset and with values chosen so that $0\le f_*+h\le1$.
Thus $f_*$ is the unique maximizer with minimal support.

\smallskip\noindent\emph{Proof.}
For $u\in N$ the characters inflated from $G/N$ take their full degree
there. Subtracting $f(1)=0$ and taking real parts gives
\[
  f(u)=\sum_{\chi\ \mathrm{new}}(-a_\chi)
                 \bigl(\chi(1)-\operatorname{Re}\chi(u)\bigr).
\]
Each bracket is nonnegative. Equality in a bracket holds exactly when
the unitary matrix representing $u$ is the identity, by summing the real
parts of its eigenvalues. This proves the zero-set statement.
The regular representation is faithful, so all new irreducibles
together separate the identity from every other element of $N$.
The mean of a new character over $N$ is zero, since averaging its
representation over $N$ projects to its $N$-invariants, a proper
$G$-submodule and hence zero. The squared degrees of the new
irreducibles sum to $g-g/n$. We obtain
\[
 t(f)(g-g/n)
 \ \le\ \sum_{\chi\ \mathrm{new}} (-a_\chi)\chi(1)
 \ =\ \frac1n\sum_{u\in N} f(u)\ \le\ 1-1/n.
\]
As $a_\chi(f_*)=-\chi(1)/g$ for each new $\chi$, both estimates are
sharp. Equality in the first requires every new coefficient to have
this value. The difference $f-f_*$ then inflates from the quotient,
and $f(1)=f_*(1)=0$ forces the asserted vanishing on $N$. The converse
follows from the same coefficient calculation.\hfill$\square$

A consequence in the opposite sign convention is worth noting. If
$f(1)=0$, $f\ge0$, and all new coefficients are nonnegative, the same
identity, with its signs reversed, forces all of them to vanish
(average over $N$). No strictly positive response of that orientation
is possible under these conditions.

\subsection*{Exactly when one nonidentity rational Frobenius type suffices}
Now take $N=D=[G,G]\ne1$. The known quotient is the abelianization.
A rational conjugacy class here means the generators of all conjugates
of one cyclic subgroup. Then a test supported on one nonidentity
rational conjugacy class has strictly negative coefficients on all
nonlinear irreducibles if and only if $D\setminus\{1\}$ is one rational
conjugacy class. The parallel statement holds for ordinary conjugacy
classes. Indeed, a class containing any member of $D$ is contained in
$D$ in either sense, and the support theorem forces it to contain
$D\setminus\{1\}$. The converse uses $f_*$.

There is a structural characterization. In the rational case $D$ must be
the additive group of $\mathbb F_{p^d}$ for some prime $p$, and, under
a suitable identification, the conjugation image $H$ of $G$ on $D$
must consist of field multiplications and satisfy
$H\mathbb F_p^\times=\mathbb F_{p^d}^{\times}$. In the ordinary case
the condition is $H=\mathbb F_{p^d}^{\times}$.
To prove necessity, all nonidentity elements of $D$ have the same
order, necessarily a prime $p$ by taking powers. Thus $D$ is a
$p$-group, and its nontrivial center is characteristic, invariant
under conjugation and taking unit powers. Transitivity forces
$D=Z(D)$, hence $D$ is an elementary abelian vector space $V$.
The conjugation image $H$ is abelian because $G/D$ is abelian.
The group generated by $H$ and the prime-field scalar maps acts
transitively on $V\setminus0$, so $V$ is irreducible as an
$\mathbb F_p[H]$-module. The commutative algebra generated by $H$
in $\operatorname{End}_{\mathbb F_p}(V)$ is a finite field: every
nonzero element has invariant kernel and is invertible, and its inverse
belongs to the algebra by finiteness. Irreducibility makes $V$ one
dimensional over this field. The transitivity statements now give the
claimed scalar conditions, which also prove sufficiency.

For $G=U_3(\mathbb F_p)$, the upper unitriangular $3$ by $3$ matrices,
$D$ is its center, the additive group supported in the $(1,3)$ entry.
Indeed $[(a,b,c),(a',b',c')]$ has central coordinate $ab'-a'b$ in
coordinates $(a,b,c)=(u_{12},u_{23},u_{13})$. The nonzero center is
one rational class, consisting of $p-1$ ordinary classes. Its density
is $(p-1)/p^3$. For $p>2$, observing one ordinary class cannot meet
the uniform sign and coverage requirements of the theorem.

\subsection*{The arithmetic interpretation needs only Dedekind zeta functions}
Let $L/\mathbb Q$ be finite Galois with group $G$, let $F=L^N$, and set
$a=[F:\mathbb Q]=g/n$. Assume $\zeta_F$ has no zeros in
$\operatorname{Re}s>1/2$. For $N=D$ this is supplied by the
Dirichlet theorem of the specified repository and the abelian
factorization of $\zeta_F$. The event $f_*(\operatorname{Frob}_p)=1$
means that the unramified prime $p$ splits completely in $F$ and fails
to split completely in $L$. Its density is $\delta=1/a-1/g$. Put
\[
  R(x)=\sum_{\substack{p\le x,\ p\ \mathrm{unramified}\\
                          f_*(\operatorname{Frob}_p)=1}}
       (\log p)(p^{-1}-x^{-1})-\delta\log x-c.
\]
Here $c$ is the real constant specified below. Then the following
conditions are equivalent:
\begin{enumerate}
\item Every nontrivial zero of $\zeta_L$ has real part $1/2$.
\item For each $0<\epsilon<1/2$, there is $C_\epsilon$ such that
$R(x)\ge-C_\epsilon x^{-1/2+\epsilon}$ for all sufficiently large $x$.
\end{enumerate}

\noindent\emph{Proof.}
Let $B(s)$ be the series over the same primes with summands
$(\log p)p^{-s}$. In $\operatorname{Re}s>1$ the unramified
prime-power coefficients of $-\zeta_M'/\zeta_M$ for a Galois field $M$
are $[M:\mathbb Q](\log p)\mathbf 1_{\operatorname{Frob}_p^k|_M=1}$
at $p^k$. Thus
\[
 B(s)=-\frac{1}{a}\frac{\zeta'_F}{\zeta_F}(s)
             +\frac{1}{g}\frac{\zeta'_L}{\zeta_L}(s)+H(s),
\]
where $H$ is holomorphic for $\operatorname{Re}s>1/2$. It removes
higher prime powers by absolute convergence and removes the finitely
many ramified Euler factors. If $b$ is the constant coefficient of
$B(s)-\delta/(s-1)$ at $1$, set $c=b-\delta$. Termwise integration for
$\operatorname{Re}s>1$ gives
\[
 \int_1^\infty R(x)x^{-s}\,dx=
 \frac{B(s)}{s(s-1)}-\frac{\delta}{(s-1)^2}-\frac{c}{s-1}.
\]
Our choice of $c$ removes the pole at $1$. In $1/2<\operatorname{Re}s<1$
the only remaining poles come from zeros $\rho$ of $\zeta_L$, with
residue $\operatorname{ord}_\rho\zeta_L/(g\rho(\rho-1))$.
In particular each real pole has negative residue.

Assume the lower bound and set $\beta=1/2+\epsilon$. The function
$V(x)=R(x)+C_\epsilon x^{\beta-1}$ is eventually nonnegative.
Changing it on a bounded interval adds an entire function to its
transform. If the abscissa of convergence of this nonnegative transform
were strictly larger than $\beta$, Landau's theorem would force a
singularity at that real abscissa. There is no such singularity at or
above $1$. Below $1$ a real pole would have negative residue, impossible
for a nonnegative integral approached from the right. Consequently
the integral is holomorphic in $\operatorname{Re}s>\beta$, and the
displayed residue forbids any $\zeta_L$ zero there. Letting $\epsilon$
decrease to zero and using the functional equation gives (1).

Conversely, the prime ideal theorem under RH for each of $F,L$ gives
$\psi_M(x)=x+O_{M,\eta}(x^{1/2+\eta})$ for every $\eta>0$.
Removing the finitely many ramified primes and higher powers
costs $O_L(\sqrt{x}\log^2(2x))$. Therefore the sum $P(x)$ of $\log p$
over the selected primes is $\delta x+O_{L,\eta}(x^{1/2+\eta})$.
Partial summation identifies the weighted sum with
$\int_1^xP(t)t^{-2}\,dt$. Its remainder after subtracting
$\delta\log x$ tends to a constant with error
$O_{L,\eta}(x^{-1/2+\eta})$ for $0<\eta<1/2$. The Mellin formula
at $1$ identifies this constant with $c$ (multiply the transform
by $s-1$ and let $s\downarrow1$).\hfill$\square$

\paragraph{An explicit quartic-nonresidue test.}
Take $L=\mathbb Q(i,\sqrt[4]{2})$. The splitting field has degree $8$:
$X^4-2$ is Eisenstein, and the real field $\mathbb Q(\sqrt[4]{2})$
does not contain $i$. A generator that multiplies $\sqrt[4]{2}$ by
$i$ while fixing $i$, and complex conjugation, give the dihedral group
of order $8$. The commutator has order $2$ and fixes
$F=\mathbb Q(i,\sqrt2)=\mathbb Q(\zeta_8)$.
The event for an odd prime is exactly
\[
 p\equiv1\pmod8,\qquad 2^{(p-1)/4}\equiv-1\pmod p.
\]
Indeed the first condition gives splitting in $F$. For these primes
$2$ is a square modulo $p$ and all fourth roots of unity are available,
so that failure of complete splitting in $L$ is precisely the quartic
nonresidue condition. The first examples are $17,41,97,137,193$.
The density and the optimal worst coefficient are both $1/8$ here.
An estimate from below for the displayed smoothed count, with main
term $(\log x)/8$ and its own constant, would settle RH for this
nonabelian splitting field. Establishing that lower estimate is the
unresolved analytic step. Neither the support theorem nor a finite
prime count supplies it.

\paragraph{Related literature.}
Suzuki studies one-sided criteria for RH and for Dirichlet GRH in
\emph{On variants of Chebyshev's conjecture},
\emph{Ramanujan Journal} 68 (2025), article 95
(\url{https://arxiv.org/abs/2411.07436}).
Bella\"iche studies character coefficients in effective Chebotarev
bounds (\url{https://arxiv.org/abs/1308.1022}).
Lewis, Morotti, Pacifici, Sanus and Tong-Viet treat ordinary conjugacy
classes equal to a difference of normal subgroups
(\url{https://arxiv.org/abs/2603.25407}).
The finite extremum above addresses a specified normalized coefficient
and a specified sign cone. Its proof is elementary character theory.
We make no claim of historical priority for the criterion or its group
theory.

\section{Scope and further questions}
The global five-term classification, the infinite exact family and the density benchmark use classical complex multiplication and local analysis. The uniform near-degree lower bound uses Dirichlet GRH for small inert primes, class numbers and coefficient heights. The Frobenius optimization uses the one-sided residue argument of the Lean proof. Its application to a nonabelian splitting field leaves an analytic inequality open. The archive includes source identities, finite witnesses, replay code and the recorded output of the checks.

\begin{enumerate}
\item Does every singular modulus of degree $h\ge6$ require at least \emph{seven} terms? This is the next support threshold after Theorem~\ref{thm:four}.
\item Can one bound $s(\Delta)$ below by a fixed positive fraction of $h(\Delta)$ for every order? The current argument loses a factor of order $p$ when the conjugates occupy different supersingular residues.
\item Can one determine the true relative density of Gaussian prime conductors with $s=h+1$? The value in Section~\ref{sec:family} counts a sufficient factorial inequality. Determinants can remain nonzero even when two coefficients have the same valuation.
\item For the dihedral field in Section~\ref{sec:transfer}, can a comparison of arithmetic quantities give the required one-sided quartic-nonresidue estimate? The normalized character coefficient is already optimal within the class of tests analyzed here.
\end{enumerate}

\appendix
\section{Replay of the earlier four-term reduction}\label{sec:earlier}
The argument below records the earlier route in \cite{prior} to the 79 discriminants treated in Section~\ref{sec:finite}. The global five-term proof in Sections~\ref{sec:global-five} and \ref{sec:four-sieves} works independently of this analytic cutoff.
\paragraph{Four different sizes.}
Suppose four nonzero conjugates have increasing absolute values $r_i$ with $r_{i+1}/r_i\ge16$. For exponents $e_1<\cdots<e_4$, multiply column $j$ of $(z_i^{e_j})$ by a positive $c_j$ with $c_{j+1}/c_j=(r_jr_{j+1})^{-(e_{j+1}-e_j)/2}$. On either side of the diagonal the absolute entries decrease by factors of at least four as one moves away from it. The sum of the other entries in each row is at most $2/3$ of the diagonal entry in absolute value. Thus the matrix is nonsingular. Any relation with fewer terms could be padded by unused exponents. For tighter gaps the earlier checker instead sums all 23 nondiagonal permutation magnitudes relative to the diagonal. The sum is largest at unit exponent gaps, by rearrangement, and a bound strictly below one suffices.

\paragraph{The analytic cutoff.}
Put $D=|\Delta|$ and $R=\sqrt D$. A split prime $p\nmid D$, $p\le R/2$, gives a primitive reduced form with leading coefficient $p$. Indeed one chooses $b^2\equiv\Delta\pmod{4p}$, $|b|\le p$, and puts $c=(b^2+D)/(4p)\ge p$. The absolute value of its singular modulus is within $2079$ of $\exp(\pi R/p)$ by the cusp estimate in \cite{riffaut}. The principal form gives leading coefficient 1. For $D\ge10^{11}$ three eligible primes occur, respectively, in $(0,.046R]$, $(.048R,.189R]$, $(.230R,.500R]$. To record the constants, set
\[
 \begin{split}
 E(y)&=\sqrt y\big[(\log y/(2\pi)+2)\log D+2(\log^2 y/(8\pi)+2)\big],\\
 B(y)&=2\sqrt y\log y+\log D+2\log 2,\qquad E(0)=0 .
 \end{split}
\]
Greni\'e and Molteni's prime ideal estimate for $K$ gives $|\psi_K(y)-y|\le E(y)$ for $y\ge100$. Its hypothesis follows from Dirichlet GRH. Over rational primes $p\le\sqrt y$, prime ideal powers contribute at most $2\sqrt y\log y$. The ramified primes and the split primes dividing $f$ cost at most $\log|d_K|+2\log f=\log D$ at larger primes. The extra term allows removal of $2$. Thus twice the eligible prime log-weight on $(x,y]$ is at least $y-x-E(x)-E(y)-B(y)$. At the three intervals the certified lower margins are respectively $68,191,751$ at $D=10^{11}$. After division by $R$ the error terms decrease: each monomial in $E(cR)/R$ contains $R^{-1/2}$ and at most two logarithm factors, each at least 9 for the positive endpoints. The same monotonicity holds for $B(cR)/R$. For two consecutive prime intervals with adjacent endpoints $vR,u'R$, the cusp ratio is bounded below by $(e^{\pi/v}-2079)/(e^{\pi/u'}+2079)>16$; the principal-to-first ratio is even larger since $p\ge3$. This proves the large range.\cite{gm,prior}

\paragraph{The finite range.}
For $10^7\le D<10^{11}$, enumerate entire residue classes where $1,p,p^2,p^3$ give four eligible reduced forms for a split prime $p$ with $4p^6\le D$. Include the powers of 2 when $D\equiv7\pmod 8$. Lifting solutions of $b^2\equiv\Delta\pmod{4p^j}$ produces the forms. For consecutive leading coefficients $a<b$, the tests $D\ge16a^2$ and $9D(b-a)^2\ge16a^2b^2$ imply a ratio exceeding 16 by the cusp bound. They hold for these powers in the specified range. The rest are tested individually. The CRT sieve accounts for all $49{,}995{,}000{,}000$ discriminants in the half-open range, including $47{,}116{,}196$ settled with three individually tested split primes. Each bin balances against the exact count in residue classes 0 and 3 modulo 4.

Below $10^7$, the integer form sieve and interval determinant test leave 80 inputs of degree at least four. The reduction enumerates primitive forms with $0\le b\le a\le c$ for absolute values, treating nonreal conjugate pairs once. The published output also covers the endpoint $10^7$. There are 67 low-degree inputs, $4{,}999{,}603$ exclusions at the first pass and another 250 by rigorous intervals. At $\Delta=-4075$ a direct complex determinant removes the sole remaining sextic. In degree six there are three smaller conjugates with moduli in $(2365,2366)$, $(2846,2847)$ and $(3701,3703)$, together with a principal conjugate larger than $10^{80}$ in modulus. If either of the first two exponent gaps is at least 12, diagonal dominance of the determinant expansion suffices. For first gaps 1 through 11, dividing by the principal root's highest power reduces to 121 nonzero complex minors, with a uniformly smaller error for every last gap. The small sieve, large sieve, boundary numbers, and all 121 minors were rerun. The archive includes the original checkers and expected outputs.\cite{prior}


\begin{thebibliography}{99}\small\raggedright
\bibitem{repo} L. Sernau and collaborators, \emph{GRH for Dirichlet $L$-functions, in Lean 4}. \href{https://github.com/lsernau/grh-lean/tree/72fead102f663d5c7a407964a13dc0cf9d977207}{Source revision}, \href{https://github.com/lsernau/grh-lean/blob/72fead102f663d5c7a407964a13dc0cf9d977207/grh/src/Statements.lean}{public types}, \href{https://github.com/lsernau/grh-lean/blob/72fead102f663d5c7a407964a13dc0cf9d977207/grh/src/New/Reduction/Zeros/Criterion.lean}{one-sided criterion}, \href{https://github.com/lsernau/grh-lean/blob/72fead102f663d5c7a407964a13dc0cf9d977207/grh/GUIDE.md}{reading guide}.
\bibitem{mitcheck} A. Sutherland, \emph{Verification files for the shorter Lean proof of RH and GRH}, reproductions and comparison logs for both editions. \url{https://math.mit.edu/~drew/shortgrh/}.
\bibitem{prior} V. Reshetnikov, ProveIt, \href{https://github.com/VladimirReshetnikov/ProveIt/blob/c477d08ef7be4241e94bd1cffb43ff56240321e3/docs/incoming/grh_cm_prime_applications.zip}{prior research archive}, especially \texttt{riffaut.tex}, \texttt{cm-finite.tex}, \texttt{cm\_small.tex}, \texttt{cm4075.tex}, \texttt{cm-conclusion.tex}. Source identifiers and a comparison of earlier archives are recorded in \texttt{checks/prior\_audit.md}.
\bibitem{companion} Companion report, \emph{Sparse equations for singular moduli and the sign of weighted prime sums}, 25 November 2026, \texttt{grh\_cm\_four\_terms\_and\_prime\_bias.zip}. See its quintic section and further research questions 1 and 2. The file identity is in \texttt{checks/prior\_audit.md}.
\bibitem{sparseparallel} Parallel report, \emph{Sparse equations for singular moduli}, 25 November 2026, \texttt{grh\_cm\_sparse\_equations.zip}. See its five-term reduction, verification through width $65{,}536$, and extremal family. File identity appears in \texttt{checks/prior\_audit.md}.
\bibitem{riffaut} Y. Bilu, F. Luca and A. Pizarro-Madariaga, \emph{Trinomials, singular moduli and Riffaut's conjecture}, Annali Sc. Norm. Sup. Pisa (5) 23 (2022), 2003--2048. \href{https://arxiv.org/abs/2003.06547}{arXiv:2003.06547}.
\bibitem{aven} M. Avenda\~no and T. Krick, \emph{Sharp bounds for the number of roots of univariate fewnomials}, J. Number Theory 131 (2011), 1209--1228, esp. Prop. 3.8 and the discussion of Lenstra's local bound. \href{https://arxiv.org/abs/1008.4808}{arXiv:1008.4808}.
\bibitem{cmred} J. Belding, R. Br\"oker, A. Enge and K. Lauter, \emph{Computing Hilbert class polynomials}, esp. \S3 on inert primes and CM reduction. \href{https://arxiv.org/abs/0802.0979}{arXiv:0802.0979}. See also D. Cox, \emph{Primes of the Form $x^2+ny^2$}, 2nd ed., on ring class fields, class numbers and conjugates.
\bibitem{supersing} J. Silverman, \emph{The Arithmetic of Elliptic Curves}, 2nd ed., Springer (2009), Chapter V, Theorem 4.1, supersingular invariants.
\bibitem{suth} A. V. Sutherland, \emph{Computing Hilbert class polynomials with the Chinese remainder theorem}, Math. Comp. 80 (2011), 501--538, Appendix 1, Lemma 8. \href{https://arxiv.org/abs/0903.2785}{arXiv:0903.2785}.
\bibitem{watkins} M. Watkins, \emph{Class numbers of imaginary quadratic fields}, Math. Comp. 73 (2004), 907--938, Table 4. \href{https://doi.org/10.1090/S0025-5718-03-01517-5}{Publisher's version}.
\bibitem{gm} L. Greni\'e and G. Molteni, \emph{Explicit versions of the prime ideal theorem for Dedekind zeta functions under GRH}, Math. Comp. 85 (2016), 889--906, Corollary 1.2. \href{https://arxiv.org/abs/1312.4463}{arXiv:1312.4463}.
\bibitem{gms} L. Greni\'e and G. Molteni, \emph{Explicit smoothed prime ideals theorems under GRH}, Theorem 1.1. \href{https://arxiv.org/abs/1312.4465}{arXiv:1312.4465}.
\bibitem{pollack} P. Pollack, \emph{Bounds for the first several prime character nonresidues}, Theorem 1.1. \href{https://arxiv.org/abs/1508.05035}{arXiv:1508.05035}.
\end{thebibliography}

\end{document}
